% =========================================================
% Algebraic Laws of Set Operations
% =========================================================

\section{Algebraic Laws of Set Operations}

% ---------------------------------------------------------
% TOOLKIT
% ---------------------------------------------------------

\begin{remark*}[Algebraic structure of set operations]
The operations $\cup$, $\cap$, and complement satisfy algebraic laws analogous
to those of logical connectives. Together they endow $\mathcal{P}(U)$ with
the structure of a \emph{Boolean algebra}. These laws justify manipulation
of set expressions in later proofs, particularly those involving equivalence
classes, partitions, and functions.
\end{remark*}

\begin{theorem}[Commutativity of Union and Intersection]\label{thm:commutativity}
Let $A,B$ be sets. Then
\[
A \cup B = B \cup A
\quad\text{and}\quad
A \cap B = B \cap A.
\]
\hyperref[prf:commutativity]{Go to proof.}
\end{theorem}

\NoLocalDependencies

\begin{remark*}[Standard quantified statement]
$\forall A, B:\; \forall x\,(x \in A \cup B \leftrightarrow x \in B \cup A)$ and $\forall x\,(x \in A \cap B \leftrightarrow x \in B \cap A)$.
\end{remark*}

\begin{theorem}[Associativity of Union and Intersection]\label{thm:associativity}
Let $A,B,C$ be sets. Then
\[
(A \cup B) \cup C = A \cup (B \cup C)
\quad\text{and}\quad
(A \cap B) \cap C = A \cap (B \cap C).
\]
\hyperref[prf:associativity]{Go to proof.}
\end{theorem}

\NoLocalDependencies

\begin{remark*}[Standard quantified statement]
$\forall A, B, C:\; \forall x\,\bigl((x \in (A \cup B) \cup C) \leftrightarrow (x \in A \cup (B \cup C))\bigr)$, and analogously for $\cap$.
\end{remark*}

\begin{theorem}[Distributive Laws]\label{thm:distributivity}
Let $A,B,C$ be sets. Then
\[
A \cap (B \cup C) = (A \cap B) \cup (A \cap C),
\]
\[
A \cup (B \cap C) = (A \cup B) \cap (A \cup C).
\]
\hyperref[prf:distributivity]{Go to proof.}
\end{theorem}

\NoLocalDependencies

\begin{remark*}[Standard quantified statement]
$\forall A, B, C:\; \forall x\,\bigl(x \in A \cap (B \cup C) \leftrightarrow x \in (A \cap B) \cup (A \cap C)\bigr)$, and dually for $\cup$ over $\cap$.
\end{remark*}

\begin{theorem}[Identity and Absorption Laws]\label{thm:identity-absorption}
Let $A,B$ be sets and $U$ a universe with $A \subseteq U$. Then
\[
A \cup \varnothing = A,
\qquad
A \cap U = A,
\]
\[
A \cup (A \cap B) = A,
\qquad
A \cap (A \cup B) = A.
\]
\hyperref[prf:identity-absorption]{Go to proof.}
\end{theorem}

\NoLocalDependencies

\begin{remark*}[Standard quantified statement]
$\forall A \subseteq U:\; \forall x\,(x \in A \cup \varnothing \leftrightarrow x \in A)$ and $\forall x\,(x \in A \cap U \leftrightarrow x \in A)$.
Absorption: $\forall A, B:\; \forall x\,(x \in A \cup (A \cap B) \leftrightarrow x \in A)$, and dually.
\end{remark*}

\begin{theorem}[Involution of Complement]\label{thm:involution}
For any $A \subseteq U$,
\[
(A^c)^c = A.
\]
\hyperref[prf:involution]{Go to proof.}
\end{theorem}

\NoLocalDependencies

\begin{remark*}[Standard quantified statement]
$\forall A \subseteq U:\; \forall x\,(x \in (A^c)^c \leftrightarrow x \in A)$.
\end{remark*}

\begin{remark*}[Non-commutative and non-associative operations]
Set difference $\setminus$ is neither commutative nor associative:
\[
A \setminus B \neq B \setminus A,
\qquad
(A \setminus B) \setminus C \neq A \setminus (B \setminus C) \quad \text{in general}.
\]
It interacts with union and intersection via:
\[
A \setminus (B \cup C) = (A \setminus B) \cap (A \setminus C),
\qquad
A \setminus (B \cap C) = (A \setminus B) \cup (A \setminus C).
\]
These follow from $A \setminus B = A \cap B^c$ together with De Morgan's laws.
\end{remark*}

\begin{remark*}[Cartesian product]
The Cartesian product $\times$ is not commutative, is associative only up to
canonical isomorphism, and distributes over union in each coordinate:
\[
A \times (B \cup C) = (A \times B) \cup (A \times C).
\]
\end{remark*}
